\documentclass[12pt]{article}%
\usepackage{amsfonts}
\usepackage{fancyhdr}
\usepackage{comment}
\usepackage[a4paper, top=2.5cm, bottom=2.5cm, left=2.2cm, right=2.2cm]%
{geometry}
\usepackage{times}
\usepackage{amsmath}
\usepackage{changepage}
\usepackage{amssymb}
\usepackage{graphicx}

\setcounter{MaxMatrixCols}{30}
\newtheorem{theorem}{Theorem}
\newtheorem{acknowledgement}[theorem]{Acknowledgement}
\newtheorem{algorithm}[theorem]{Algorithm}
\newtheorem{axiom}{Axiom}
\newtheorem{case}[theorem]{Case}
\newtheorem{claim}[theorem]{Claim}
\newtheorem{conclusion}[theorem]{Conclusion}
\newtheorem{condition}[theorem]{Condition}
\newtheorem{conjecture}[theorem]{Conjecture}
\newtheorem{corollary}[theorem]{Corollary}
\newtheorem{criterion}[theorem]{Criterion}
\newtheorem{definition}[theorem]{Definition}
\newtheorem{example}[theorem]{Example}
\newtheorem{exercise}[theorem]{Exercise}
\newtheorem{lemma}[theorem]{Lemma}
\newtheorem{notation}[theorem]{Notation}
\newtheorem{problem}[theorem]{Problem}
\newtheorem{proposition}[theorem]{Proposition}
\newtheorem{remark}[theorem]{Remark}
\newtheorem{solution}[theorem]{Solution}
\newtheorem{summary}[theorem]{Summary}
\newenvironment{proof}[1][Proof]{\textbf{#1.} }{\ \rule{0.5em}{0.5em}}

\newcommand{\Q}{\mathbb{Q}}
\newcommand{\R}{\mathbb{R}}
\newcommand{\C}{\mathbb{C}}
\newcommand{\Z}{\mathbb{Z}}

%%%%
% ME %
%%%%

\usepackage{listings}
\usepackage{color}
\usepackage{xcolor}
\usepackage{mdframed}

\definecolor{dkgreen}{rgb}{0,0.6,0}
\definecolor{gray}{rgb}{0.5,0.5,0.5}
\definecolor{mauve}{rgb}{0.58,0,0.82}

\lstset{%frame=tb,
language=Matlab,
aboveskip=3mm,
belowskip=3mm,
showstringspaces=false,
columns=flexible,
basicstyle={\small\ttfamily},
numbers=none,
numberstyle=\tiny\color{gray},
keywordstyle=\color{blue},
commentstyle=\color{dkgreen},
stringstyle=\color{mauve},
breaklines=true,
breakatwhitespace=true,
tabsize=3
}

\usepackage{outlines}
\usepackage{enumitem}
%\setenumerate[1]{label=\Roman*.}
%\setenumerate[2]{label=\Alph*.}
%\setenumerate[3]{label=\roman*.}
%\setenumerate[4]{label=\alph*.}

\usepackage{titlesec} 

\titleformat{\subsection}[runin]
  {\normalfont\large\bfseries}{\thesubsection}{1em}{}
\titleformat{\subsubsection}[runin]
  {\normalfont\normalsize\bfseries}{\thesubsubsection}{1em}{}

% for enumerate spacing
\newenvironment{tight_enum}{
\begin{enumerate}[label=\alph*.]
\setlength{\itemsep}{0pt}
\setlength{\parskip}{0pt}
}{\end{enumerate}}

\usepackage{exercise}

\setlength\parindent{0pt}

\newcommand{\code}[1]{\texttt{#1}}
\newcommand{\shp}{$>>>$ }
\newcommand{\tab}{\hspace*{45pt}}

\usepackage{multicol}
\setlength{\columnsep}{0.6cm}

\begin{document}

\title{Homework 10: \\ Optimization}

\author{EEC 3150}
\date{}

\maketitle

\begin{center}
\fbox{\fbox{\parbox{5.5in}{\centering
For each Python exercise, write a separate Python script. Submit all scripts on Blackboard under the appropriate assignment. Name your script files with the following format: \\
\vspace{10pt}
HW$<$assignment number$>$\_Ex$<$exercise number$>$\_$<$FirstInitial$><$LastName$>$.py \\
\vspace{10pt}
Example: HW01\_Ex01\_GWest.py \\
\vspace{10pt}
For short answer or math exercises, either do your work on a sheet of paper and upload a picture of it to Blackboard or submit some kind of text file or Word document to Blackboard. You can place all such exercises in a single document as long as you clearly label each exercise. These files should follow a similar naming convention as the Python scripts.
}}}
\end{center}


\pagebreak


\begin{multicols*}{2}



%%% NEW EXERCISE %%%
\begin{Exercise}[origin={25 points, Python}]

Consider the function
\begin{equation*}
	f(x,y) = (x^2+y^2)^{0.5} - 0.6(\cos(x)+\cos(y))
\end{equation*}

\begin{enumerate}
	\item Implement and run a 2D gradient descent method to find the minimum of the function
	\item On a single 3D plot...
	\begin{enumerate}
		\item Plot the function in 3D over the domain $[-15,15] \times [-15,15]$ using plot\_surface
		\item Plot the path of the GD method on the $xy$ plane beneath the surface
	\end{enumerate}
\end{enumerate}

\end{Exercise}

%%% NEW EXERCISE %%%
\begin{Exercise}[origin={25 points, Python}]

Redo the previous exercise with 2D simulated annealing. You will need to specify the covariance matrix of the multivariate normal distribution. Use something like \\

\code{width**2*np.eye((2,2))} \\

where \code{width} is the proposal width used in the same sense as 1D simulated annealing.

\end{Exercise}


\vfill\null
\columnbreak


%%% NEW EXERCISE %%%
\begin{Exercise}[origin={25 points, Python}]

We can represent the wavefunction of the electron in a hydrogen atom as
\begin{equation*}
	\psi(r) = \dfrac{2}{R^{3/2}} \exp(-r/R)
\end{equation*}
The energy of this wavefunction is a function of the $R$ parameter
\begin{equation*}
	E(R) = 2B\bigg(\dfrac{A^2}{2R^2} - \dfrac{A}{R}\bigg)
\end{equation*}
where
\begin{align*}
	A & = \hbar^2/(k m e^2) \\
	B & = m k^2 e^3/(2\hbar^2)
\end{align*}
and
\begin{itemize}
	\item $\hbar$ (reduced Planck constant) = \\ $1.05\times10^{-34}$ J s
	\item $k_e$ (Coulomb constant) = \\ $8.98\times10^9$ N m$^2$/C$^2$
	\item $e$ (electron charge) = \\ $1.60\times10^{-19}$ C
	\item $m$ (electron mass) = \\ $9.10 \times 10^{-31}$ kg
\end{itemize}
The value of $R$ which minimizes this energy is the Bohr radius and the minimum value of $E$ is the ground state energy of the hydrogen atom.

\begin{enumerate}
	\item Implement and run simulated annealing to determine the value of the Bohr radius and the ground state energy
	\item Print the values
	\item TIP: the Bohr radius is approximately $10^{-10}$ m so be aware of this when setting the value of your proposal width
\end{enumerate}

\end{Exercise}






%\pagebreak
%\columnbreak


\end{multicols*}

\end{document}






